Theorem~\ref{thm:asymmetry} leaves only frames with no nontrivial projective isometries. In particular, a construction that forces a nontrivial cyclic permutation or a conjugation symmetry cannot produce these parameters. The root-of-unity obstruction applies even when a finite-phase ansatz does not force a permutation symmetry.

For a Hermitian matrix with unit diagonal and exact equiangular modulus, the numerical search minimized the tightness residual $\|G^2-\lambda G\|_F$. In the free $C_3$ class, 48 starts per projection algorithm and local refinement gave a best residual of about $0.7168843846$. With 28 independent phases after anchoring, 1{,}000 unrestricted starts and refinement gave residuals in $[0.7168843846433283,0.7168843846433312]$. Each stored endpoint matched the best free-$C_3$ point to a maximum entrywise error below $7\times10^{-8}$ under a search allowing relabeling, rephasing, and complex conjugation.

A Gram matrix satisfying \eqref{eq:gram} would give a construction. A proof of nonexistence by moment infeasibility would require further constraints, for example additional higher-degree equation multiples or PSD constraints. Existence of a complex $\operatorname{ETF}(4,9)$, equivalently of its Naimark complement $\operatorname{ETF}(5,9)$, remains open to our knowledge \cite{fickusmixon}.

\paragraph{Formalization.}
The companion Lean 4 draft contains declarations for the signature change of variables and for the rational identity in \eqref{eq:fixed1identity} in finite real coordinates. The proposed real refutation starts from the lifted equations and a normalized kernel vector. The derivation of those equations from an antiunitary, the remaining symmetry exclusions, and the archived moment verifiers are outside the scope of the draft. Its pinned build instructions are in \pathcode{etf_search/lean/README.md}. The draft awaits a Lean build. The exact replays reported in this article use the independent Python verifiers.
